% page 417 of the facsimile (image p-416 of the scan)
% block 162.6 x 237.7 mm ; left margin 39.4 mm
\pgc{17.780mm}{0.000mm}{
\l{\cel{54}409.}
\l{}
\l{\sou{ordi}nar.\cel{2}(\sou{trans.}) Dispozar - spacale, tempale o hierarkiale -}
\l{\cel{3}segun ula regulo, ula normo, la kozi - ici relate iti.- DEFIS}
\l{}
\l{\sou{ordinacar}. (\sou{trans}.) (\sou{liturgio katol.}). Elevar a sacerdoteso,}
\l{\cel{3}diakoneso, e c., per konferar (a lu) ta ordo sakra. - DEFIS.}
\l{}
\l{\sou{ordinara.}\cel{2}Qua esas konforma a la ordineso kustumala ; qua esas}
\l{\cel{3}ye la grado kustumala. - EFIS.}
\l{}
\l{\sou{ordonanco}. (\sou{armeo)}.I.Servisto soldata di oficiro. II.\cel{1}\sou{Ordonanc-}}
\l{\cel{3}\sou{oficiro}\cel{1}:\cel{2}kavaliero qua portas ad oficiro superiora la ko-}
\l{\cel{3}mandi de alta oficiero di la stato. - DFIRS.}
\l{}
\l{\sou{oreado}. (\sou{mitol.}) Nimfo di la monti, di la boski. DEFS.}
\l{}
\l{\sou{orelo}. (\sou{anat.)}\cel{2}I. La organo di audado. - II. Parto extera di}
\l{\cel{3}la aud-organo, ye singla latero di la kapo. - EFIS.}
\l{}
\l{\sou{oreliono}. (\sou{patol.)}\cel{2}Influro di la celulo-tisuo qua envelopas la}
\l{\cel{3}glando parotida. - F.}
\l{}
\l{\sou{\textquotedbl{}oremus\textquotedbl{}}. (\sou{liturgio katol.}) Prego.}
\l{}
\l{\sou{orfano.}\cel{2}Infanto o puero di qua la patri mortis. -EFIS.}
\l{}
\l{organo. - I. En korpo vivoza : grupo de elementi, koordinita}
\l{\cel{3}tale ke exekutesas funciono specala. - II. Anke pri mashino.}
\l{\cel{3}- III. (\sou{metaf.})\cel{2}To quo expresas la penso, la opiniono, la}
\l{\cel{3}sentimenti di altru. - DEFIRS.}
\l{}
\l{\sou{organdio}.- Muslino apretizita. DEFIS.}
\l{}
\l{\sou{organicismo}.\cel{2}(\sou{filoz.}) Doktrino qua admisas la ideo ke vivo}
\l{\cel{3}rezultas, ne de forco qua \textquotedbl{}anmizas\textquotedbl{} la organi, ma de la}
\l{\cel{3}organi ipsa. - DEFIRS.}
\l{}
\l{\sou{organika}. Qua relatas la enti organizita. -\cel{1}\sou{Kemio organika}\cel{1}:}
\l{\cel{3}qua studias la produkturi di la enti organizita. DEFIRS.}
\l{}
\l{\sou{organismo}. I. Ensemblo de la organi ye qui konsistas ento}
\l{\cel{3}vivoza. - II.\cel{1}\sou{anke metaf}.\cel{2}- DEFIRS.}
\l{}
\l{\sou{organzino}. Silko-filo qua pasis dufoye en la mashino qua tordas}
\l{\cel{3}la silko kruda, e qua destinesas a formacar la warpo di}
\l{\cel{3}stofo. - DEFI.}
\l{}
\l{\sou{orgeno}. (\sou{muzi}ko).Sufl-instrumento muzikala, ek tubi di dimen-}
\l{\cel{3}sioni diversa, quan onu resonigas, mediace klavaro, per}
\l{\cel{3}enduktar en olu la aero quan furnisas suflilo. - EFIRS.}
\l{}
\l{\sou{orgio}. I. (en\cel{1}\sou{la epoki antiqua}). Festo honorize di Bako}
\l{\cel{3}(Bacchus). - II. Krapulo pro troa manjo e troa drinko.}
\l{\cel{3}DEFIRS.}
}
